\section{Parte 2: Historias de usuario}

Se redactaron las historias de usuario usando el formato: \textit{Como [usuario], quiero [funcionalidad], para [beneficio]}.

\begin{table}[htbp]
\centering
\small
\renewcommand{\arraystretch}{1.15}
\begin{tabularx}{\textwidth}{c p{2.2cm} X X}
\toprule
\textbf{ID} & \textbf{Rol} & \textbf{Quiero...} & \textbf{Para...} \\
\midrule
HU-01 & Participante & inscribirme a una capacitación usando mi DNI & no tener que ir presencialmente a anotarme. \\
\midrule
HU-02 & Participante & que me llegue un correo cuando me inscriba & saber que mi registro se hizo bien. \\
\midrule
HU-03 & Docente & pasar asistencia escaneando un código QR & no perder tiempo con listas en papel. \\
\midrule
HU-04 & Docente & ingresar las notas de los participantes & que el sistema diga quién aprobó y quién no. \\
\midrule
HU-05 & Participante & descargar mi certificado en PDF con QR & tener un documento válido que pueda mostrar. \\
\midrule
HU-06 & Admin & crear una capacitación nueva con fechas y cupos & que los participantes se puedan inscribir. \\
\midrule
HU-07 & Admin & registrar diferentes instituciones en el sistema & que cada municipalidad o empresa maneje sus cursos aparte. \\
\midrule
HU-08 & Ciudadano & escanear el QR de un certificado para verificarlo & comprobar que no es falso. \\
\midrule
HU-09 & Gerente & sacar reportes de las capacitaciones realizadas & presentarlos en auditorías. \\
\midrule
HU-10 & Admin & asignar roles a cada usuario (admin, docente, alumno) & que cada uno vea solo lo que le corresponde. \\
\midrule
HU-11 & Docente & marcar asistencia aunque no haya internet & que se sincronice después automáticamente. \\
\midrule
HU-12 & Participante & ver los cursos disponibles en un catálogo & decidir a cuál inscribirme. \\
\bottomrule
\end{tabularx}
\caption{Historias de usuario del SIGC-CUSCO.}
\label{tab:historias}
\end{table}
